% =====================================================================
%  Fine-tuning beyond CHSH: minimum measurement dependence for CGLMP,
%  Mermin, and contextuality inequalities
%  Companion to "Exact minimum measurement dependence for violating the
%  CHSH inequality" (paper/main.tex, in preparation).
%  Target: Phys. Rev. A (revtex4-2), same class options as main.tex.
%  Compile: ~/bin/tectonic paper/companion.tex
%  (bibtex pulls only the cited entries from the shared paper/refs.bib)
% =====================================================================
\documentclass[aps,pra,reprint,superscriptaddress,floatfix]{revtex4-2}

\usepackage{amsmath,amssymb}
\usepackage{amsthm}
\usepackage{booktabs}
\usepackage{array}
\usepackage{graphicx}

\newcommand{\TV}{\mathrm{TV}}
\newcommand{\ket}[1]{\lvert #1\rangle}
\newtheorem{lemma}{Lemma}
\newtheorem{theorem}{Theorem}
\newtheorem{proposition}{Proposition}
\newtheorem{conjecture}{Conjecture}

\begin{document}

% TODO: replace with real author block (name, affiliation, email) before submission.
\title{Fine-tuning beyond {CHSH}: minimum measurement dependence for {CGLMP}, {Mermin}, and contextuality inequalities}
\author{Dan Petersen}
\noaffiliation

\date{\today}

\begin{abstract}
We compute the exact minimum measurement dependence, in the worst-case
total-variation metric $F=\max_{s}\,\TV(p(\cdot\mid s)\,\Vert\,\rho)$ of the
companion {CHSH} paper, that local or noncontextual hidden-variable models
need to reach quantum values in three further scenarios: the {CGLMP}
inequality for $d=3$, the {Mermin} (GHZ-parity) inequalities for $n=3$ and
$n=5$, and the Peres--Mermin square. For {CGLMP} $d=3$ the {CHSH} form
survives verbatim: for any target $T$ on the first linear segment,
$F^{*}=(T-2)/(2\min(N,4))$ for every $N$ --- $\sqrt3/9-1/12\approx0.109$
($N\ge4$) at the maximally-entangled-state value
$T_{\mathrm{ME}}=(12+8\sqrt3)/9$, and $(\sqrt{11/3}-1)/8\approx0.114$ at
the quantum maximum $T_{\mathrm{Q}}=1+\sqrt{11/3}$ of Ac\'in \emph{et al.}
Where the quantum value equals the algebraic maximum ---
{Mermin} $n=3$ and the Peres--Mermin square --- the price is the cap,
$\lceil N/4\rceil/N$ (saturating at $1/4$) and $\lceil N/6\rceil/N$
(saturating at $1/6$); the latter is, to the coverage of our search, the
first exact minimum-measurement-dependence value for a contextuality
scenario. For {Mermin} $n=5$, frustration --- every hidden state misses at
least six of the sixteen conditions --- changes the answer: the exact values
at $N=2,\dots,6$ are $\tfrac12,\tfrac23,\tfrac12,\tfrac35,\tfrac12$, the
$N=5$ value lies strictly above the counting bound $\lceil 3N/8\rceil/N$,
and $F^{*}\to\tfrac38=m_{\min}/K$ as $N\to\infty$. We state a $K/m_{\min}$
principle as a conjecture whose instances are proved, and compare with
Alai's faithful-model program for GHZ-{Mermin} correlations, whose metric
is incommensurable with ours but whose conclusions align. The results are
exact and certificate-backed; their significance is structural: they
delimit when the {CHSH} price-per-violation form generalizes and when
frustration takes over.
\end{abstract}

\pacs{03.65.Ta, 03.65.Ud, 03.67.-a}
\keywords{Bell inequality, CGLMP inequality, Mermin inequality,
contextuality, Peres--Mermin square, measurement dependence, fine-tuning}

\maketitle

% =====================================================================
\section{Introduction}
\label{sec:intro}

The companion paper to this work
\cite{mainpaper2026} quantified the measurement dependence a local
hidden-variable model must have to violate the {CHSH} inequality: it proved
$S\le 2+8F$ for every local {CHSH} model with
$F=\max_{ab}\TV(p(\cdot\mid a,b)\Vert\rho)$ the worst-case
total-variation (TV) distance of the conditional source from a fixed
baseline $\rho$, computed the exact maximum $S_{\max}(N,F)$ for $N$-state
uniform-source models, and showed that the minimum fine-tuning for the
Tsirelson value is $F^{*}(2\sqrt{2};N)=(\sqrt{2}-1)/\min(N,4)$ --- flat at
$(\sqrt{2}-1)/4$ for $N\ge 4$, so that the bottleneck is the four {CHSH}
correlators, not the size of the hidden source. The question that paper
left open was whether that ``four correlators'' saturation is special to
{CHSH} or is an instance of a general principle for inequalities whose
algebraic maximum requires satisfying $K$ independent conditions with a
parity-type obstruction (every hidden state misses at least one).

This paper answers that question for three {Bell} inequalities and one
contextuality scenario, in the same metric and with the same two-lemma
template (a parity/obstruction lemma counting missed conditions, the row
lemma of \cite{mainpaper2026} bounding the per-row gain from fine-tuning,
and a round-robin construction over near-miss patterns):

\begin{itemize}
\item \textbf{{CGLMP}, $d=3$} (Sec.~\ref{sec:cglmp}): the generalization
  holds \emph{verbatim}. Both the maximally-entangled-state value
  $T_{\mathrm{ME}}=(12+8\sqrt{3})/9\approx 2.87293$ \cite{cglmp2002} and the
  quantum maximum $T_{\mathrm{Q}}=1+\sqrt{11/3}\approx2.91485$
  \cite{adgl2002,npa2008} lie below the algebraic maximum $4$, on the first linear
  segment, and $F^{*}(T;N)=(T-2)/(2\min(N,4))$ exactly, for every $N$.
\item \textbf{{Mermin} (GHZ-parity), $n=3$} (Sec.~\ref{sec:mermin3}): the
  quantum value \emph{is} the algebraic maximum ($4$ on
  $\ket{\mathrm{GHZ}_{3}}$, perfect correlators), so the target sits at the
  cap rather than on the linear segment; the price is the cap
  $\lceil N/4\rceil/N$, identical to {CHSH}'s with $K=4$.
\item \textbf{{Mermin}, $n=5$} (Sec.~\ref{sec:mermin5}): frustration ---
  every hidden state misses at least six of the sixteen conditions
  ($m_{\min}=6$) --- breaks the all-but-one structure; the exact values at
  $N=2,\dots,6$ are $\tfrac12,\tfrac23,\tfrac12,\tfrac35,\tfrac12$ (the
  $N=5$ value strictly above the general bound $\lceil 3N/8\rceil/N=2/5$),
  with $\lim_{N\to\infty}F^{*}=3/8=m_{\min}/K$.
\item \textbf{Peres--Mermin square} (Sec.~\ref{sec:pm}): the state-
  independent contextuality analog has exactly the same shape --- a parity
  lemma, the same row lemma, a bound $G\le 5/6+F$, and a cap
  $F^{*}_{\mathrm{PM}}(N)=\lceil N/6\rceil/N$ saturating at $1/6$ with the
  six lines playing the role of the four {CHSH} correlators.
\end{itemize}

Together these establish the generalization in one direction ({CGLMP}) and
an identifiable failure mode in the other (GHZ-{Mermin}, $n\ge 5$), and they
delimit the scope of the fine-tuning theory: the structural data are the
triple $(K,m_{\min},\text{balance})$, not $K$ alone
(Sec.~\ref{sec:unify}). A search of the {arXiv} literature found no prior
computation of the minimum fine-tuning in the per-setting TV metric used
here for any of these scenarios (the GHZ-{Mermin} family has been treated
in a different, incommensurable metric in \cite{alai2026b},
Sec.~\ref{sec:alai}); that statement, with its coverage limits, is a search
result, not a proven novelty claim (Sec.~\ref{sec:scope}).

\emph{Scope note.} We do not treat the {I3322} inequality: its quantum
target $0.250875384514\ldots$ is a \emph{conjectured} maximum, attained by
an infinite-dimensional family of states with only a numerical upper bound
on optimality, and its single-party marginal terms add biasable rows that
complicate the clean $K$-row structure studied here. {CGLMP} $d\ge 4$ and
{Mermin} $n\ge 7$ are likewise out of reach or out of scope
(Sec.~\ref{sec:scope}).

% =====================================================================
\section{Setup: the metric, generalized}
\label{sec:setup}

\subsection{Models and cost}

Fix a scenario with parties, settings, and outcomes as in the scenarios
below, and an inequality whose value is a linear combination of the
$K$ setting blocks $s$ that enter it. A finite model consists of:

\begin{itemize}
\item $N$ hidden states $\lambda=1,\dots,N$ with \emph{uniform} source
  $\rho(\lambda)=1/N$;
\item for each setting block $s$, a conditional source $p[s,\lambda]\ge 0$
  with $\sum_{\lambda}p[s,\lambda]=1$ (the rows $p[s,\cdot]$ are
  independent --- no joint constraint across blocks);
\item deterministic responses (per party, per local setting).
\end{itemize}

The fine-tuning cost is the exact generalization of the project's metric:

\begin{equation}
F=\max_{s}\ \TV\big(p(\cdot\mid s)\,\Vert\,\rho\big),\qquad
\TV=\frac12\sum_{\lambda}\big|p[s,\lambda]-\tfrac1N\big|.
\label{eq:Fgen}
\end{equation}

Rows that do not enter the inequality are never biased at an optimum
(biasing them costs $F$ without helping the value), so only the $K$ rows
matter. The model-class caveats of \cite{mainpaper2026} carry over: the
lower-bound side does not depend on the source distribution, and
stochastic responses reduce to deterministic ones by the vertex argument.

\subsection{Generalized row lemma}

\begin{lemma}[row optimum for weights in $\{-1,0,+1\}$]
\label{lem:row}
Let $w_{1},\dots,w_{N}\in\{-1,0,+1\}$ be the weight multiset of a row,
$\mu=\frac1N\sum_{\lambda}w(\lambda)$ its mean, $w_{\max}=\max_{\lambda}
w(\lambda)$, and $w^{\uparrow}_{1}\le\cdots\le w^{\uparrow}_{N}$ the
ascending sort. The maximum of $\sum_{\lambda}p[\lambda]\,w(\lambda)$ over
all $p$ with $\TV(p\Vert\rho)\le F$, at $F=k/N$, is
\begin{equation}
\begin{aligned}
h(k/N)=\;&\mu+\frac{k\,w_{\max}-\sum_{i=1}^{k}w^{\uparrow}_{i}}{N},
\quad 0\le k\le N-1,\\
&\qquad h(1)=w_{\max}.
\end{aligned}
\label{eq:rowopt}
\end{equation}
For $w\in\{\pm 1\}$ this reduces to the row lemma of
\cite{mainpaper2026}, $h=\mu+2\min(F,n/N)$ with $n$ the wrong-sign count.
\end{lemma}

\begin{proof}
Write $d=p-\rho$, so $d$ is zero-sum with $\|d\|_{1}\le 2F$. Decompose
$d=d^{+}-d^{-}$ with $d^{\pm}\ge 0$ and $\sum d^{+}=\sum d^{-}=F$. Then
\begin{equation*}
d\!\cdot\! w=\sum d^{+}w-\sum d^{-}w
\ \le\ F\,w_{\max}-\min_{d^{-}}\sum d^{-}w .
\end{equation*}
The minimum of $\sum d^{-}w$ over $0\le d^{-}\le 1/N$,
$\sum d^{-}=F=k/N$ is attained by draining the smallest-$w$ entries first
(greedy; each has capacity $1/N$), giving $\sum_{i=1}^{k}w^{\uparrow}_{i}/N$.
Achievability: take $d^{+}=F\,e_{\arg\max}$ (capacity $(N-1)/N\ge k/N$) and
$d^{-}$ as in the greedy; then $p=\rho+d^{+}-d^{-}\ge 0$ by construction.
\end{proof}

\subsection{Finite-max formula and uniform bound}

For fixed responses the inequality value is linear in each row
$p[s,\cdot]$ and the feasible set is a product of TV balls, so the rows
decouple; a hidden state's relevant data is its \emph{pattern} --- the
vector of its per-row contributions. Then, exactly as in Theorem~1 of
\cite{mainpaper2026},

\begin{equation}
V_{\max}(N,F)=\max_{Q}\Big[\,V_{0}(Q)+\sum_{r}h_{r}(Q,F)\,\Big],
\label{eq:finmax}
\end{equation}

where $Q$ ranges over multisets of $N$ realizable patterns,
$V_{0}(Q)$ is the uniform-source value, and each $h_{r}$ is given by
Lemma~\ref{lem:row} (or its $\{\pm1\}$ special case). For each fixed $Q$
the bracket is an exact piecewise-linear function of $F$ with breakpoints at
multiples of $1/N$; the envelope over $Q$ is convex on each cell
$[k/N,(k+1)/N]$ but may bend \emph{inside} a cell where two multisets cross
(this happens for {Mermin} $n=5$ at $N=2$, Sec.~\ref{sec:mermin5}). Full
multiset enumeration \emph{is} the exact computation; the grid values
$V(k/N)$, $k=0,\dots,N$, determine the whole curve exactly when a single
construction matches the envelope at every grid point (the convexity
argument of Sec.~\ref{sec:pm}). (No LP is used; the
enumerations run in exact $\mathbb{Q}$ or quadratic-field arithmetic,
Sec.~\ref{sec:scope}.)

\begin{proposition}[uniform global bound]
\label{prop:uniform}
For every multiset $Q$ and every $F$,
\begin{equation}
V_{\max}(N,F)\ \le\ B+2KF,
\label{eq:uniform}
\end{equation}
where $B=\max$ pattern value (the local bound).
\end{proposition}

\begin{proof}
Per row, with $d=p-\rho$ and $\|d\|_{1}\le 2F$,
$|E_{r}-\mu_{r}|=|d\cdot w|\le F\,(w_{\max}-w_{\min})\le 2F$; summing over
the $K$ rows and using $V_{0}(Q)\le B$ (every pattern's value is at most
$B$) gives the claim.
\end{proof}

Instances: {CHSH} (recovered from \cite{mainpaper2026}), {CGLMP} $d=3$ and
{Mermin} $n=3$ all give $2+8F$; {Mermin} $n=5$ gives $4+32F$ (the $N=5$
curve of Sec.~\ref{sec:mermin5} hits it exactly at $F=1/5$).

\emph{Structural parameters.} Two structural parameters of an inequality,
plus a balance condition, govern the answers below: $K$ = the number of
extremal (row) conditions that must all be satisfied to reach the
algebraic maximum $A$; $m_{\min}$ = the minimum number of conditions any
single pattern misses (the \emph{frustration index}); and \emph{balance}
= whether the min-miss patterns can be distributed evenly across
conditions. {CHSH}, {CGLMP} $d=3$ and {Mermin} $n=3$ have $m_{\min}=1$
(``all-but-one'' near-miss structure); {Mermin} $n=5$ has $m_{\min}=6$.

% =====================================================================
\section{{CGLMP}, $d=3$: the headline form generalizes verbatim}
\label{sec:cglmp}

\subsection{The inequality and its constants}

Two parties, two settings each ($A_{1},A_{2}$ / $B_{1},B_{2}$), three
outcomes. The {CGLMP} inequality of \cite{cglmp2002} at $d=3$ (all
equalities mod $3$) is

\begin{align}
I_{3}={}&
\big[P(A_{1}\!=\!B_{1})-P(A_{1}\!=\!B_{1}\!-\!1)\big]\notag\\
&+\big[P(B_{1}\!=\!A_{2}\!+\!1)-P(B_{1}\!=\!A_{2})\big]\notag\\
&+\big[P(A_{2}\!=\!B_{2})-P(A_{2}\!=\!B_{2}\!-\!1)\big]\notag\\
&+\big[P(B_{2}\!=\!A_{1})-P(B_{2}\!=\!A_{1}\!-\!1)\big].
\label{eq:i3}
\end{align}

\begin{itemize}
\item \textbf{Local bound $I_{3}\le 2$.} Stated in \cite{cglmp2002}
  (``$I_{3}\le 2$'' for local realism) and independently re-proved here by
  exhaustive check of all $81$ deterministic strategies ($27$ patterns),
  certificate \texttt{ineq\_cglmp.py} Part~A, cross-checked by direct
  brute force in \texttt{ineq\_xcheck.py}.
\item \textbf{Quantum values.} \citet{cglmp2002} give
  $I_{3}=4/(-9+6\sqrt{3})=(12+8\sqrt{3})/9\approx 2.87293=:T_{\mathrm{ME}}$
  for the maximally entangled two-qutrit state with Fourier-based
  measurements, and state explicitly that they have no proof that these
  measurements are optimal. Indeed this is not the maximum: Ac\'in, Durt,
  Gisin and Latorre \cite{adgl2002} showed that the non-maximally entangled
  state $(\ket{00}+\gamma\ket{11}+\ket{22})/\sqrt{2+\gamma^{2}}$ with
  $\gamma=(\sqrt{11}-\sqrt{3})/2\approx0.7923$ and the same measurements
  reaches $T_{\mathrm{Q}}=1+\sqrt{11/3}\approx2.91485$. This is the quantum
  maximum of $I_{3}$ to numerical precision: the semidefinite upper bound of
  the Navascu\'es--Pironio--Ac\'in hierarchy at level $1{+}AB$ equals
  $2.9149$ with the rank-loop stopping criterion satisfied \cite{npa2008},
  although no closed-form optimality proof is known. Both values were
  recomputed numerically for this paper; we treat both as targets
  (Sec.~\ref{sec:scope}).
\item \textbf{Algebraic maximum $4$.} Stated in \cite{cglmp2002};
  re-proved unattainable locally here by the modular obstruction below.
\end{itemize}

\subsection{Pattern structure: the modular obstruction ($K=4$,
$m_{\min}=1$)}

Write $d_{00}=A_{1}-B_{1}$, $d_{01}=A_{1}-B_{2}$, $d_{10}=A_{2}-B_{1}$,
$d_{11}=A_{2}-B_{2}$ (mod $3$). Every deterministic strategy's differences
satisfy
\begin{equation}
d_{00}+d_{11}\equiv d_{01}+d_{10}\pmod{3},
\label{eq:mod}
\end{equation}
the additive analog of {CHSH}'s parity lemma ($27$ valid patterns out of
$81$). The four row weight vectors, in $\{-1,0,+1\}$, are
\begin{equation*}
w_{00}(d)=+1[d{=}0]-1[d{=}2],\quad w_{01}(d)=+1[d{=}0]-1[d{=}1],
\end{equation*}
\begin{equation*}
w_{10}(d)=+1[d{=}2]-1[d{=}0],\quad w_{11}(d)=+1[d{=}0]-1[d{=}2].
\end{equation*}
The algebraic-maximum conditions are $d_{00}=0$, $d_{01}=0$, $d_{10}=2$,
$d_{11}=0$; constraint \eqref{eq:mod} gives $0+0\equiv 0+2\pmod 3$, false
--- \emph{no pattern attains all four row conditions}. Enumeration finds
exactly \textbf{four near-miss patterns, one per blocked row}, each with
value $2=B$:
\begin{equation*}
\begin{array}{ccl}
(d_{00},d_{01},d_{10},d_{11}) & \text{blocks row} & \text{value}\\[2pt]
(2,0,2,0) & A_{1},B_{1} & 2\\
(0,1,2,0) & A_{1},B_{2} & 2\\
(0,0,0,0) & A_{2},B_{1} & 2\\
(0,0,2,2) & A_{2},B_{2} & 2.
\end{array}
\end{equation*}
(Six further patterns, of weight type $(+1,+1,0,0)$, also attain the local
bound $2$, but they miss two conditions each with weight $0$ rather than
$-1$: such rows can be bent only at rate $1$ per unit TV, and the exhaustive
enumeration confirms they never beat the near-miss round-robin.) This is the
direct {CGLMP} analog of {CHSH}'s four miss-types (single cover, like
{CHSH}); in each near-miss pattern the blocked row has weight
$-1$ and the other three rows $+1$, so under round-robin every row
contains both signs (slope $2$ per row on the first segment).

\subsection{Exact results}

\begin{theorem}[minimum fine-tuning at the quantum value]
\label{thm:cglmp}
Let $T\in\{T_{\mathrm{ME}},T_{\mathrm{Q}}\}$ with
$T_{\mathrm{ME}}=(12+8\sqrt{3})/9\approx 2.87293$ and
$T_{\mathrm{Q}}=1+\sqrt{11/3}\approx2.91485$. Then, for the $N$-state
uniform-source class,
\begin{equation}
F^{*}(I_{3}=T;\,N)=\frac{T-2}{2\,\min(N,4)}
\label{eq:fcglmp}
\end{equation}
for every $N$. At $T_{\mathrm{ME}}$: $F^{*}=\sqrt{3}/9-1/12\approx
0.109117$ exactly for $N\ge 4$, $\approx0.21823$ at $N=2$ and
$\approx0.14549$ at $N=3$. At $T_{\mathrm{Q}}$:
$F^{*}=(\sqrt{11/3}-1)/8\approx0.114357$ for $N\ge4$, $\approx0.22871$ at
$N=2$ and $\approx0.15248$ at $N=3$.
\end{theorem}

\begin{proof}
The uniform bound \eqref{eq:uniform} gives $F^{*}\ge (T-2)/8$ for all
$N$. For $N\ge 4$ the round-robin construction over the four near-misses
(state $i$ carries the near-miss pattern of row $i\bmod 4$) gives
$V_{\max}(N,F)\ge 2+8F$ on $[0,\lfloor N/4\rfloor/N]$ (every row contains
both signs, Sec.~\ref{sec:cglmp}); the target crosses there because
$(T-2)/8<1/7\le \lfloor N/4\rfloor/N$ for every $N\ge 4$ and both targets
($0.1092$ and $0.1144$ versus $1/7\approx0.1429$; asserted in
$\mathbb{Q}(\sqrt{3})$ for $T_{\mathrm{ME}}$ at $N=4,5,6$, and by direct
comparison for $T_{\mathrm{Q}}$; the minimum of $\lfloor N/4\rfloor/N$ for
$N\ge 4$ is $1/7$). For $N=2,3$ the value $(T-2)/(2N)$ follows from the
exact per-multiset segment analysis in $\mathbb{Q}(\sqrt{3})$ over
\emph{all} $378$ / $3654$ multisets (certificate \texttt{ineq\_cglmp.py}
Part~B, for $T_{\mathrm{ME}}$): the minimum is attained by a round-robin
near-miss multiset, and no other multiset crosses $T$ earlier. For
$T_{\mathrm{Q}}$ the same per-multiset analysis was repeated in floating
point over all multisets for $N=2,3,4$ (\texttt{cglmp\_Tprime.py}) and
reproduces the closed form to $10^{-9}$ --- a numerical cross-check, not an
exact certificate.
\end{proof}

So the {CHSH} $\min(N,K)$ headline form holds \emph{exactly}, with $K=4$,
because here the quantum-side value $T$ is below the algebraic maximum
(the target lies in the linear regime, as for {CHSH}). Note
$F^{*}({\rm CGLMP},N\ge 4)\approx 0.1091$ (at $T_{\mathrm{ME}}$) or
$0.1144$ (at $T_{\mathrm{Q}}$) $>F^{*}({\rm CHSH},N\ge 4)\approx 0.10355$:
the {CGLMP} quantum points sit slightly farther along the same slope-$8$
segment.

\begin{proposition}[cap at the algebraic maximum]
\label{prop:cglmpcap}
The minimum fine-tuning to reach $I_{3}=4$ is $\lceil N/4\rceil/N$ for all
$N$.
\end{proposition}

\begin{proof}
Write $n_{r}^{+}=\#\{\lambda: w_{r}(\lambda)<+1\}$ (states blocking row
$r$ from its per-row maximum). Every pattern blocks at least one row, so
$\sum_{r}n_{r}^{+}\ge N$ and $\max_{r}n_{r}^{+}\ge \lceil N/4\rceil$;
reaching $I_{3}=4$ forces each row supported on $\{w_{r}=+1\}$, hence
$\TV_{r}\ge n_{r}^{+}/N$ and $F^{*}\ge \lceil N/4\rceil/N$. Conversely the
round-robin over the four near-misses has $\max_{r}n_{r}^{+}=\lceil
N/4\rceil$ (asserted for $N=2,\dots,64$; the counting argument is valid
for all $N$), and supporting each row on its $+1$ states costs exactly
$n_{r}^{+}/N$ TV. The exhaustive minimum over all multisets equals
$\lceil N/4\rceil/N$ for $N=2,\dots,6$.
\end{proof}

The exact curves $V_{\max}(N,k/N)$, $N=2,\dots,6$ (full multiset
enumeration, every breakpoint; certificate \texttt{ineq\_cglmp.out}), with
first segment $V_{\max}(N,F)=2+2\min(N,4)F$ on $[0,1/N]$:

\begin{equation*}
\begin{array}{c|cccc}
N & V(0) & V(1/N) & V(2/N) & F^{*}(T;N)\\
\hline
2 & 2 & 4 & --- & (T-2)/4\\
3 & 2 & 4 & --- & (T-2)/6\\
4 & 2 & 4 & --- & (T-2)/8\\
5 & 2 & 18/5 & 4 & (T-2)/8\\
6 & 2 & 10/3 & 4 & (T-2)/8.
\end{array}
\end{equation*}

% =====================================================================
\section{{Mermin} (GHZ-parity), $n=3$: the cap form}
\label{sec:mermin3}

\subsection{The inequality and its constants}

Three parties, two settings each ($X$ or $Y$), binary outcomes. The
tripartite {Mermin} (GHZ-parity) inequality, in its standard even-weight
form, is
\begin{equation}
M_{3}=E(XXX)-E(XYY)-E(YXY)-E(YYX),
\label{eq:m3}
\end{equation}
the four terms being the four tripartite perfect correlators of the
$\ket{\mathrm{GHZ}_{3}}$ state, as catalogued for this inequality in
\cite{alai2026b} (after Hall, Phys.\ Rev.\ A 84, 022102 (2011), as cited
there). Its constants are re-proved here by exact computation, not taken
on authority:

\begin{itemize}
\item \textbf{Local bound $B_{3}=2$.} Proved by exhaustive
  deterministic-strategy check (certificate \texttt{ineq\_mermin.py}
  Part~A; max pattern value $=2$) and independently by direct brute force
  over all $64$ strategies with no pattern parameterization
  (\texttt{ineq\_xcheck.py}). It equals the standard $2^{(n-1)/2}$.
\item \textbf{Quantum value $4=$ algebraic maximum.} Proved by exact
  $\mathbb{Q}(i)$ matrix computation (\texttt{ineq\_xcheck.py} Part~(2);
  explicit Kronecker products of Pauli matrices, no trigonometric
  shortcut): on $\ket{\mathrm{GHZ}_{3}}=(\ket{000}+\ket{111})/\sqrt{2}$
  the four correlators are \emph{perfect} --- $E(XXX)=+1$ and
  $E(XYY)=E(YXY)=E(YYX)=-1$ --- so $M_{3}=4$. Since four terms of
  $|\cdot|\le 1$ cannot exceed $4$, \emph{the quantum value equals the
  algebraic maximum}: the key structural difference from {CHSH}, and the
  reason the headline changes form below.
\end{itemize}

\subsection{Pattern structure: the parity obstruction ($K=4$,
$m_{\min}=1$)}

A hidden state's responses $(o_{i}[X],o_{i}[Y])$, $i=1,\dots,3$, give $64$
deterministic strategies. Its correlator vector over the four even-weight
settings is $e_{s}=P\cdot\prod_{i\in S_{Y}}r_{i}$ with
$(P,r_{1},r_{2},r_{3})\in\{\pm1\}^{4}$; the $16$ parameter pairs carry a
$2$-fold redundancy ($(P,r)\sim(P,-r)$, invisible on even-weight rows),
so there are exactly \textbf{$8$ distinct realizable patterns} --- the
other $8$ of the $16$ possible $\{\pm1\}^{4}$ vectors violate the parity
product constraint
\begin{equation*}
e_{XXX}\,e_{XYY}\,e_{YXY}\,e_{YYX}=+1
\end{equation*}
(each factor appears squared across the product). Each realizable pattern
is realized by exactly $8$ of the $64$ strategies. The algebraic-maximum
conditions $E(XXX)=+1$, $E(XYY)=E(YXY)=E(YYX)=-1$ have target product
$(+1)(-1)^{3}=-1$ while every pattern has product $+1$ --- the
\emph{parity obstruction} (exact analog of the obstruction lemma of
\cite{mainpaper2026}). Enumeration gives: minimum miss count $=1$; exactly
\textbf{four near-miss patterns, one per condition} ($32$ near-miss
strategies, $8$ per condition); local bound $B_{3}=\max$ pattern value
$=2$. This is the direct {Mermin} analog of {CHSH}'s four miss-types --- a
single cover, exactly like {CHSH} and {CGLMP} $d=3$. Indeed the pattern
data ($8$ even-parity patterns with values $\pm2$, four single-miss types)
are combinatorially identical to {CHSH}'s, so every finite-$N$ statement of
\cite{mainpaper2026} transfers to $M_{3}$ by relabelling; the results below
make that transfer explicit.

\subsection{Exact results}

\begin{theorem}[minimum fine-tuning at the algebraic maximum]
\label{thm:mermin3}
For the $N$-state uniform-source class,
\begin{equation}
F^{*}(M_{3}=4;\,N)=\left\lceil\frac{N}{4}\right\rceil\!\Big/N
\label{eq:fmermin3}
\end{equation}
for all $N\ge 2$: values $\tfrac12,\tfrac13,\tfrac14,\tfrac25,\tfrac13,
\tfrac27,\tfrac14,\dots$, identical to the {CHSH} cap of
\cite{mainpaper2026} with $K=4$.
\end{theorem}

\begin{proof}
\emph{Lower bound.} Every pattern misses at least one condition, so
$\sum_{s}n_{s}(Q)\ge N$ and $\max_{s}n_{s}\ge\lceil N/4\rceil$;
reaching $M_{3}=4$ forces each row fully bent (row value $+1$ requires
$p$ supported on the correct-sign states, $\TV_{r}\ge n_{r}/N$).
\emph{Upper bound.} Round-robin over the four near-miss patterns (state
$i$ misses condition $i\bmod 4$) gives $\max_{s}n_{s}\le\lceil N/4
\rceil$ and $\sum_{s}n_{s}=N$ (asserted for $N=2,\dots,64$; the counting
argument is valid for all $N$), so
$V_{\max}(N,\lceil N/4\rceil/N)\ge 2+2(N/N)=4$ exactly.
\end{proof}

The exact curves, $N=2,\dots,6$ (certificate \texttt{ineq\_mermin.out};
the first-segment bend values coincide with {CHSH}'s $N=5,6$ bends ---
same $K$, same $m_{\min}$):
\begin{equation*}
\begin{array}{c|cccc}
N & V(0) & V(1/N) & V(2/N) & F^{*}(4;N)\\
\hline
2 & 2 & 4 & --- & 1/2\\
3 & 2 & 4 & --- & 1/3\\
4 & 2 & 4 & --- & 1/4\\
5 & 2 & 18/5 & 4 & 2/5\\
6 & 2 & 10/3 & 4 & 1/3.
\end{array}
\end{equation*}

\emph{Intermediate targets.} For any $T\in(2,4)$ lying on the linear
segment (i.e.\ $(T-2)/8\le \lfloor N/4\rfloor/N$),
$F^{*}(T;N)=(T-2)/(2\min(N,4))$ --- in particular
$F^{*}(2\sqrt{2};N)=(\sqrt{2}-1)/4$ for all $N\ge 4$, \emph{the same
number as {CHSH}}. So the fine-tuning price per unit of violation
generalizes with $K=4$; only the algebraic-maximum target departs from the
$\min(N,K)$ form.

\emph{Why the $\min(N,K)$ headline does not survive here.} {CHSH}'s
quantum point $2\sqrt{2}<4$ lies on the first linear segment (slope
$2K=8$ once $N\ge K$), so $F^{*}=(T-B)/(2\min(N,K))$. For {Mermin} $n=3$
the quantum point \emph{is} the algebraic maximum: reaching it requires
\emph{every} row to be fully bent (not merely active), which costs the cap
$\lceil N/K\rceil/N$ --- a quantity that oscillates
($\tfrac14,\tfrac25,\tfrac13,\tfrac27,\dots$) rather than saturating at
$N=K$. The condition structure ($K=4$, $m_{\min}=1$, all-but-one
patterns, parity obstruction) is the same; what differs is where the
target sits.

% =====================================================================
\section{{Mermin}, $n=5$: frustration governs the price}
\label{sec:mermin5}

\subsection{The inequality and its constants}

Five parties, two settings each, binary outcomes:
$M_{5}=\sum_{|S|\ \mathrm{even}}(-1)^{|S|/2}E(S)$ over the $16$ even-
weight setting strings. As for $n=3$, the constants are re-proved here by
exact computation:

\begin{itemize}
\item \textbf{Local bound $B_{5}=4$.} Proved by exhaustive check ($32$
  patterns; \texttt{ineq\_xcheck.py}: $1024$ strategies).
\item \textbf{Quantum value $16=$ algebraic maximum.} Proved by exact
  $\mathbb{Q}(i)$ computation: on $\ket{\mathrm{GHZ}_{5}}$ every
  even-weight correlator is exactly $(-1)^{|S|/2}$, so $M_{5}=16$.
\end{itemize}

\subsection{Pattern structure: frustration ($K=16$, $m_{\min}=6$)}

There are $32$ distinct realizable patterns ($2^{5}$; same $(P,r)$
redundancy as Sec.~\ref{sec:mermin3}) and $16$ active rows. Enumeration
(certificate \texttt{ineq\_mermin.py} Part~D, cross-checked by
\texttt{ineq\_xcheck.py}) gives: \textbf{minimum miss count $=6$} --- the
all-but-one structure \emph{fails} ($m_{\min}>1$); exactly
\textbf{$16$ min-miss patterns}, with uniform incidence \textbf{$6$ per
condition} (each of the $16$ conditions is missed by exactly $6$ of the
$16$ min-miss patterns). The parity-type obstruction persists (no
$0$-miss pattern) but is \emph{frustrated}: every state must miss at
least six of the sixteen conditions.

\subsection{First-segment slope: $2A_{\max B}$, not $2\min(K,m_{\min}N)$}

On $[0,1/N]$ the finite-max formula \eqref{eq:finmax} is affine per
multiset $Q$ (value $V_{0}(Q)+2F\,A(Q)$, with
$A(Q)=\#\{s:0<n_{s}<N\}$ the number of active rows), so the right
derivative at $F=0$ is $2\,A_{\max B}$ where
\begin{equation*}
\begin{aligned}
A_{\max B}=\;&\max\{\,A(Q):V_{0}(Q)=B_{5}\,\}\\
\le\;&\min(K,m_{\min}N).
\end{aligned}
\end{equation*}
the maximum over multisets whose patterns all achieve the local bound
(each such pattern misses exactly $m_{\min}=6$, so $\sum_{s}n_{s}=6N$).
Exact enumeration of the $16$ bound-achieving patterns
(\texttt{ineq\_mermin.out} Part~E2) gives

\begin{equation*}
\begin{array}{c|ccccc}
N & 2 & 3 & 4 & 5 & 6\\
\hline
A_{\max B} & 8 & 12 & 16 & 16 & 16\\
\min(K,m_{\min}N) & 12 & 16 & 16 & 16 & 16.
\end{array}
\end{equation*}

The naive formula $2\min(K,m_{\min}N)$ is \emph{false} at $N=2,3$ (it
would give slopes $24,32$; the true first-segment slopes are $16,24$). At
$N=2$ the envelope bends \emph{inside} $[0,1/2]$: no bound-achieving $Q$
can support the chord to $V(1/2)=16$ (that would need $A=12>
A_{\max B}=8$), so a sub-bound multiset takes over before $F=1/2$. For
$N=3,4,5,6$ the bound-achieving line $B_{5}+2A_{\max B}F$ reaches
$V(1/N)$ exactly (convexity of the envelope then forces it to be the whole
segment): single segments of slope $24$ ($N=3$) and $32=2K$ ($N=4,5,6$;
$V(1/5)=52/5=4+32/5$ hits the global bound $4+32F$ exactly). For $N\ge 4$
the miss sets of a few bound-achieving patterns already cover all $16$
rows ($m_{\min}N\ge 16$), so every row is active from $F=0$. This refines
--- and at small $N$ corrects --- the ``$2\min(N,K)$'' rule of
\cite{mainpaper2026}, which assumed $m_{\min}=1$.

\subsection{Exact results}

The exact curves, $N=2,\dots,6$ (full multiset enumeration;
\texttt{ineq\_mermin.out} and \texttt{ineq\_mermin5\_n5.out}):
\begin{equation*}
\begin{array}{c|ccccc}
N & V(0) & V(1/N) & V(2/N) & V(3/N) & F^{*}(16;N)\\
\hline
2 & 4 & 16 & --- & --- & 1/2\\
3 & 4 & 12 & 16 & --- & 2/3\\
4 & 4 & 12 & 16 & --- & 1/2\\
5 & 4 & 52/5 & 72/5 & 16 & \mathbf{3/5}\\
6 & 4 & 28/3 & 14 & 16 & 1/2.
\end{array}
\end{equation*}

\begin{theorem}[minimum fine-tuning for {Mermin} $n=5$]
\label{thm:mermin5}
$F^{*}(M_{5}=16;N)=\min_{Q}\max_{s}n_{s}(Q)/N$, with exact values
$F^{*}=\tfrac12,\tfrac23,\tfrac12,\tfrac35,\tfrac12$ at $N=2,3,4,5,6$.
The $N=5$ value $\tfrac35$ lies \emph{strictly above} the general lower
bound $\lceil 3N/8\rceil/N=\tfrac25$; equality holds at $N=2,3,4,6$. For
general $N=16m+r$,
\begin{equation}
\frac{\lceil 3N/8\rceil}{N}\ \le\ F^{*}(M_{5}=16;N)\ \le\
\frac{6m+r}{16m+r},
\label{eq:sandwich}
\end{equation}
with equality in the lower bound at $N=2,3,4,6$ and for all
$N\equiv 0,1,2,3,4,6\pmod{16}$, and
\begin{equation}
\lim_{N\to\infty}F^{*}(M_{5}=16;N)=\tfrac38=\frac{m_{\min}}{K}.
\label{eq:lim38}
\end{equation}
\end{theorem}

\begin{proof}
\emph{Lower bound, all $N$.} Every pattern misses at least $6$
conditions, so $\sum_{s}n_{s}(Q)\ge 6N$ and $\max_{s}n_{s}\ge\lceil
6N/16\rceil=\lceil 3N/8\rceil$; reaching $M_{5}=16$ forces each row fully
bent ($\TV_{s}\ge n_{s}/N$). \emph{Upper bounds.} (a) For
$N=16m+r$ with $r\in\{0,1,2,3,4,6\}$: $m$ copies of all $16$ min-miss
patterns (incidence $6m$ per condition) plus an optimal size-$r$ multiset
$Q_{r}$ give $\max_{s}n_{s}=6m+\max_{s}n_{s}(Q_{r})$, and
$\max_{s}n_{s}(Q_{r})=\lceil 3r/8\rceil$ for $r=0,1,2$ (one realizable
pattern; a disjoint complementary pair) and, by the exact enumeration of
this section, also for $r=3,4,6$ (values $2,2,3$); since
$\lceil 3(16m+r)/8\rceil=6m+\lceil 3r/8\rceil$, this is
$\lceil 3N/8\rceil$ exactly (certificate Part~F for $r=0,1,2$; the
hardcoded witness masks were verified to be realizable pattern miss masks
in \texttt{ineq\_mermin5\_n5.out} Part~B, and $16$ disjoint realizable
pairs exist, so the $r=2$ construction is not an isolated artifact).
(b) For other residues $N=16m+r$: trivially $F^{*}\le (6m+r)/(16m+r)$.
The limit follows from \eqref{eq:sandwich}; the exact $N=5,6$ values and
full curves are the exhaustive enumeration.
\end{proof}

\emph{Reading.} The fine-tuning price to reach the GHZ-$5$ algebraic
maximum is governed by frustration: each state must miss at least $6$ of
$16$ conditions, so even an infinitely large source needs $F\ge 3/8$.
This is the same qualitative lesson as Alai's universal ceiling in
\cite{alai2026b} (their $F\le 1/2$ in Hall's metric), but in our TV metric
with an exact asymptotic value: $\lim F^{*}=m_{\min}/K=3/8$ at $n=5$, vs
the oscillating cap $\lceil N/4\rceil/N$ (liminf $1/4$) at $n=3$.

% =====================================================================
\section{The Peres--Mermin square: the contextuality analog}
\label{sec:pm}

\subsection{Setup: a fine-tuning problem}

The Peres--Mermin (PM) square \cite{peres1990,mermin1990b} is a $3\times 3$ grid of
nine binary observables on two qubits whose six lines (three rows, three
columns) have operator products
\begin{equation*}
\begin{array}{c|ccc}
 & C_{0} & C_{1} & C_{2}\\
\hline
R_{0} & X\!\otimes\!X & X\!\otimes\!I & I\!\otimes\!X\\
R_{1} & Y\!\otimes\!Y & X\!\otimes\!Z & Z\!\otimes\!X\\
R_{2} & Z\!\otimes\!Z & I\!\otimes\!Z & Z\!\otimes\!I
\end{array},
\end{equation*}
with row products $+I,+I,+I$ and column products $-I,+I,+I$, the $-I$
product being on column $C_{0}$ (any relabeling of lines is equivalent;
certificate \texttt{fringe\_pm.py} Part~0b verifies all six product
identities and each line's commutativity by exact integer-matrix
arithmetic). The product of the six line products is $-I$ while the
product of the nine observables taken twice (once per row, once per
column) is $+I$: a state-independent contextuality proof.

\emph{The game.} A referee picks one of the six lines $c$ uniformly; the
player wins iff the product of the three preassigned outcomes on that
line equals the target $t_{c}$, with $(t_{0},\dots,t_{5})=(1,1,1,-1,1,1)$
($-1$ on $C_{0}$, the line whose product is $-I$). A
\emph{noncontextual} (NC) model assigns to each hidden state $\lambda$ a
full assignment $v(\lambda)\in\{\pm1\}^{9}$; by the parity lemma below no
assignment satisfies all six lines, so $G\le 5/6$ at $F=0$. A
\emph{quantum strategy} measures the three observables of the chosen line
in any state --- each line product is literally $\pm I$ and the target
equals its sign on every line --- so \emph{quantum $G=1$} (certificate
PM-Qa--d, exact matrices).

\emph{Measurement dependence.} The source may depend on which line is
measured: $p[c,\lambda]\ge 0$, row-sum $1$, uniform baseline
$\rho(\lambda)=1/N$; cost $F=\max_{c}\TV(p(\cdot\mid c)\Vert\rho)$ --- the
same metric as everywhere else in this paper. The rows $p[c]$ are
independent, so for a fixed state multiset $Q$ each line can be optimized
separately, exactly as in the {CHSH} theory. Question: what is the minimum
$F$ for an $N$-state NC-response model to reach $G=1$, and what is the
exact curve $G_{\max}(N,F)$? Two honest differences from {CHSH}, stated up
front: (i) the quantum value equals the algebraic maximum
(state-independent contextuality), so the interesting quantity is the
\emph{cap} rather than an intermediate quantum value --- the {KCBS}
inequality would test the intermediate case (Sec.~\ref{sec:scope});
(ii) there is no even-parity restriction on states: all $512$ assignments
exist, of which exactly $32$ odd-parity sign vectors are realizable on the
six line products.

\subsection{Parity lemma and miss-set census}

\begin{lemma}[PM parity]
\label{lem:pmparity}
For any assignment $v$, write $s_{c}(v):=t_{c}\prod_{\text{cell}\in
c}v_{\text{cell}}\in\pm 1$ ($+1$ iff line $c$ is satisfied). Then every
assignment violates an \emph{odd} number of lines --- in particular at
least one.
\end{lemma}

\begin{proof}
Each cell lies in exactly one row and one column, so
\begin{equation*}
\prod_{c}s_{c}(v)=\frac{\prod_{c}\prod_{\text{cell}\in c}v_{\text{cell}}}
{\prod_{c}t_{c}}=\frac{+1}{-1}=-1.
\end{equation*}
\end{proof}

This is the contextuality analog of the obstruction lemma of
\cite{mainpaper2026}. The certificate additionally censuses all $512$
assignments: exactly the \textbf{$32$ odd-weight miss-sets} occur ($6$
single-line ``near-miss'' types, $20$ three-line, $6$ five-line), and a
near-miss type exists for \emph{every} one of the six lines (explicit
examples in the log; e.g.\ the all-$+1$ assignment violates only the
$-I$ line). This census is what makes the round-robin construction
below possible, exactly as the four {CHSH} miss-types make the round-robin
$Q_{N}$ of \cite{mainpaper2026} possible.

\subsection{Global bound and first segment}

For line $c$ write $w_{c}:=\#\{\lambda:s_{c}(\lambda)=-1\}$ and
$\mu_{c}:=(N-2w_{c})/N$. Then $P(\text{win }c)=\tfrac12(1+E_{p[c]}
[s_{c}])$, and the row lemma of \cite{mainpaper2026} (max of a linear
functional over a TV ball --- the only non-combinatorial ingredient of the
whole project, reused verbatim) gives
\begin{equation*}
\max_{\TV(p[c]\Vert\rho)\le F}P(\text{win }c)
=1-\frac{w_{c}}{N}+\min\Big(F,\frac{w_{c}}{N}\Big).
\end{equation*}

\begin{theorem}[PM global bound]
\label{thm:pm2}
For every $N$ and every $F$: $G_{\max}(N,F)\le \tfrac56+F$.
\end{theorem}

\begin{proof}
For a model with miss-count vector $(w_{0},\dots,w_{5})$ of its state
multiset $Q$ and total $M:=\sum_{c}w_{c}=\sum_{\lambda}\mathrm{misscount}
(\lambda)$, the row bound gives per line at most $1-w_{c}/N+F$;
averaging over the six lines,
\begin{equation*}
G\ \le\ 1-\frac{M}{6N}+F.
\end{equation*}
By Lemma~\ref{lem:pmparity}, $\mathrm{misscount}(\lambda)\ge 1$ for every
state, so $M\ge N$, and therefore $G\le 5/6+F$.
\end{proof}

This is the exact analog of $S\le 2+8F$: a model-independent bound whose
tightness is carried by the round-robin construction. (The {CHSH} version
has slope $8=2K$ because each of $K=4$ rows bends at rate $2$ per unit TV;
here each of $K=6$ lines bends at rate $1$ and the objective averages over
six lines, so the total slope is $1$.)

\emph{Round-robin.} Let state $i$ ($i=0,\dots,N-1$) carry the near-miss
assignment of line $i\bmod 6$. Then $w_{c}(Q_{N})\in\{\lfloor N/6\rfloor,\lceil N/6\rceil\}$ for every $c$ and $M=N$, so
\begin{equation*}
G_{Q_{N}}(F)=\frac56+\frac16\sum_{c}\min\Big(F,\frac{w_{c}}{N}\Big).
\end{equation*}

\begin{theorem}[PM first segment]
\label{thm:pm3}
For $N\ge 6$: $G_{\max}(N,F)=\tfrac56+F$ on
$[0,\lfloor N/6\rfloor/N]$ (all six lines active; lower bound by
$Q_{N}$, upper bound by Theorem~\ref{thm:pm2}). For $N=2,\dots,5$ the
round-robin has only $N$ active lines and the whole curve is the single
segment $G=\tfrac56+\tfrac{N}{6}F$ on $[0,1/N]$, reaching $1$ at $F=1/N$.
\end{theorem}

\emph{Exact curves, $N=2,\dots,10$ (certificate PM-5).} Because the
objective depends on $Q$ only through the miss-count vector $w$, the exact
curve is the upper envelope over the finitely many \emph{achievable}
$w$-vectors --- computed by an exact $\mathbb{Q}$ dynamic program over the
$32$ odd miss-set indicator vectors, then asserted to equal the
round-robin curve \emph{at every breakpoint $k/N$} ($k=0,\dots,N$) for
$N=2,\dots,10$. Grid points suffice because the envelope $E(F)=\max_{w}
G_{w}(F)$ is convex (max of affines) while the round-robin value $R(F)$ is
affine on each cell $[k/N,(k+1)/N]$: $h:=E-R$ is convex on each cell,
nonnegative everywhere ($Q_{N}$ is among the maximized), and zero at both
cell endpoints by the assertion; a convex function with $h\ge 0$ and
$h(a)=h(b)=0$ vanishes throughout. The breakpoint assertions therefore
certify the whole curve (same technique as the exact-curve certificate of
\cite{mainpaper2026}). All values are exact rationals:

\begin{equation*}
\begin{array}{c|l|c}
N & G_{\max}(N,F) & F^{*}_{\mathrm{PM}}(N)\\
\hline
2 & \tfrac56+\tfrac13 F\ \text{on }[0,\tfrac12] & \tfrac12\\
3 & \tfrac56+\tfrac12 F\ \text{on }[0,\tfrac13] & \tfrac13\\
4 & \tfrac56+\tfrac23 F\ \text{on }[0,\tfrac14] & \tfrac14\\
5 & \tfrac56+\tfrac56 F\ \text{on }[0,\tfrac15] & \tfrac15\\
6 & \tfrac56+F\ \text{on }[0,\tfrac16] & \tfrac16\\
7 & \tfrac56+F;\ \text{then } \tfrac{41}{42}+\tfrac16(F-\tfrac17) & \tfrac27\\
8 & \tfrac56+F;\ \text{then } \tfrac{23}{24}+\tfrac13(F-\tfrac18) & \tfrac14\\
9 & \tfrac56+F;\ \text{then } \tfrac{17}{18}+\tfrac12(F-\tfrac19) & \tfrac29\\
10 & \tfrac56+F;\ \text{then } \tfrac{14}{15}+\tfrac23(F-\tfrac1{10}) & \tfrac15.
\end{array}
\end{equation*}

(``then'' = on the interval to $F^{*}$; every row reaches exactly $1$ at
$F^{*}$.) For $N\ge 6$ the first segment has slope exactly $1$ --- the
direct analog of $S=2+8F$ on $[0,1/N]$; the bends for
$N\not\equiv 0\pmod 6$ are the analog of the {CHSH} bend region and are
exact here for all $N\le 10$ by full enumeration.

\subsection{The cap}

\begin{theorem}[PM cap]
\label{thm:pm4cap}
$F^{*}_{\mathrm{PM}}(N):=\min\{F:G_{\max}(N,F)=1\}=\lceil N/6
\rceil/N$ for every $N\ge 2$. In particular the cost \emph{saturates at
$1/6$}: $F^{*}_{\mathrm{PM}}(6k)=1/6$ for all $k$, and
$F^{*}_{\mathrm{PM}}(N)\to 1/6$.
\end{theorem}

\begin{proof}
\emph{Necessity.} $G=1$ forces $P(\text{win }c)=1$ for every line $c$
(each term is $\le 1$). By the row lemma this requires $p[c]$ supported on
$\{\lambda:s_{c}(\lambda)=+1\}$, which --- when $w_{c}>0$ --- requires
$F\ge w_{c}/N$. Hence $F\ge \max_{c}w_{c}(Q)/N$; and by
Lemma~\ref{lem:pmparity}, $\sum_{c}w_{c}=\sum_{\lambda}
\mathrm{misscount}(\lambda)\ge N$, so $\max_{c}w_{c}\ge\lceil N/6
\rceil$. Thus $F^{*}_{\mathrm{PM}}(N)\ge\lceil N/6\rceil/N$.
\emph{Sufficiency.} The round-robin $Q_{N}$ has $\max_{c}w_{c}=\lceil
N/6\rceil$. At $F=\lceil N/6\rceil/N$ set each row $p[c]$ uniform on its
satisfying states (there are $N-w_{c}$ of them); then
\begin{equation*}
\TV(p[c]\Vert\rho)=\tfrac12\Big[\,(N-w_{c})
\big|\tfrac{1}{N-w_{c}}-\tfrac1N\big|+w_{c}\tfrac1N\,\Big]
=\frac{w_{c}}{N}\le F,
\end{equation*}
and every line is won with probability $1$, so $G=1$.
\end{proof}

\emph{Explicit witnesses (certificate Part PM-6, $N=4,\dots,8$).} State
$i$ carries the near-miss assignment of line $i\bmod 6$; rows $p[c]$
uniform on satisfiers. Verified exactly in $\mathbb{Q}$: row sums $1$;
per-row $\TV=w_{c}/N$; max per-row $\TV=\lceil N/6\rceil/N$; $G=1$. The
$N=4$ witness is printed in full in the log --- note its elegance: the two
lines that every state satisfies ($C_{1}$, $C_{2}$) need \emph{zero}
fine-tuning (uniform rows), and each of the other four rows spends
exactly $\TV=1/4$.

\emph{Structural remark.} The proof is a line-by-line translation of the
cap theorem of \cite{mainpaper2026}: parity lemma
($\sum$ misses $\ge N$); $\min_{Q}\max_{c}w_{c}/N$; round-robin
sufficiency. The only new ingredient is the Lemma-2 census ($32$ odd
miss-sets, all six near-miss types), which replaces {CHSH}'s four
even-parity patterns. Combined with Theorem~8 of
\cite{mainpaper2026} --- the no-signalling {CHSH} cap at the algebraic
maximum --- the same statement now holds in two scenarios: \emph{hiding
the fine-tuning from the marginals never raises the price of reaching the
algebraic maximum.}

% =====================================================================
\section{Unification: the $(K,m_{\min},\text{balance})$ principle}
\label{sec:unify}

The data from the main paper \cite{mainpaper2026} and the four scenarios
of this paper ($F^{*}(A;N)$ = minimum fine-tuning to reach the algebraic
maximum $A$; status labels in the following paragraph):

\small
\begin{equation*}
\begin{array}{l|c|c|l}
\text{Scenario} & K & m_{\min} & F^{*}(A;N)\\
\hline
\text{{CHSH}} & 4 & 1 & \lceil N/4\rceil/N\to 1/4\\
\text{{CGLMP} $d=3$} & 4 & 1 & (T-2)/(2\min(N,4))\ \text{at }T<A;\\
 & & & \text{cap: }\lceil N/4\rceil/N\to 1/4\\
\text{{Mermin} $n=3$} & 4 & 1 & \lceil N/4\rceil/N\to 1/4\\
\text{PM square} & 6 & 1 & \lceil N/6\rceil/N\to 1/6\\
\text{{Mermin} $n=5$} & 16 & 6 & \tfrac12,\tfrac23,\tfrac12,\tfrac35,\tfrac12\ (N\!=\!2\dots 6);\\
 & & & \ge\lceil 3N/8\rceil/N;\ \lim F^{*}=3/8.
\end{array}
\end{equation*}
\normalsize

\begin{conjecture}[$K/m_{\min}$ principle]
\label{conj:kmm}
For a parity-frustrated extremal game with $K$ conditions, frustration
index $m_{\min}$, and a near-miss state for each condition: (i) the
minimum TV fine-tuning to reach the algebraic maximum $A$
\emph{saturates at $m_{\min}/K$} as $N\to\infty$; (ii) when $m_{\min}=1$
it equals $\lceil N/K\rceil/N$ exactly for all $N$; (iii) when the
quantum value $T<A$ lies on the first linear segment, the cost to reach
$T$ is $(T-B)/(2\min(N,K))$.
\end{conjecture}

Status labels, stated explicitly: clause (ii) is \textbf{PROVED} for
{CHSH} \cite{mainpaper2026}, the PM square (Theorem~\ref{thm:pm4cap}), and
{Mermin} $n=3$ (Theorem~\ref{thm:mermin3}); the {CGLMP} cap is the
extrapolated instance of the proved linear-regime price. Clause (iii) is
\textbf{PROVED} for {CHSH} and {CGLMP} $d=3$ (and, as an intermediate
target, for {Mermin} $n=3$, Sec.~\ref{sec:mermin3}). Clause (i) is
\textbf{PROVED} for {Mermin} $n=5$ (the limit $3/8$,
Eq.~\eqref{eq:lim38}) and for all $m_{\min}=1$ cases (which saturate at
$1/K$); the general case is \textbf{CONJECTURED}. The {Mermin} $n=5$ row
delimits the pattern precisely: the \emph{saturation value}
$m_{\min}/K=6/16=3/8$ survives as the limit, and the lower bound
$\lceil m_{\min}N/K\rceil/N$ is tight for $N=2,3,4,6$ and all
$N\equiv 0,1,2,3,4,6\pmod{16}$, but the finite-$N$ cap form can be
\emph{strictly exceeded} --- at $N=5$ the exact value $3/5$ lies above
$\lceil 3\cdot 5/8\rceil/5=2/5$; for large $m$ the remaining residues
$r\in\{5,7,8,\dots,15\}\pmod{16}$ are sandwiched but not pinned
(Sec.~\ref{sec:scope}). The structural data are therefore the triple
$(K,m_{\min},\text{balance})$, not $K$ alone.

The $m_{\min}=1$ case reduces to exactly the two-lemma template used in
both \cite{mainpaper2026} and this paper: obstruction lemma
($\sum$ misses $\ge N$) + row lemma (per-condition gain bounded by the row
optimum) + round-robin over near-miss types. For general $m_{\min}$ the
lower bound $\lceil m_{\min}N/K\rceil/N$ is a one-line extension of the
same counting argument, but the \emph{achieving construction} needs a
balanced packing of $m_{\min}$-miss states instead of round-robin --- and
{Mermin} $n=5$ shows such packings are not always available at the bound
(strictly exceeded at $N=5$). That is the main open item
(Sec.~\ref{sec:scope}).

% =====================================================================
\section{Comparison with Alai's faithful-model program}
\label{sec:alai}

\citet{alai2026b} solve the \emph{faithful} version of Hall's problem for
GHZ-{Mermin} correlations (every full correlator reproduced and every
proper-subset marginal vanishing), with metric the ``surrendered fraction''
$F:=M/2$, where $M$ is Hall's pairwise-$L^{1}$ measurement-dependence
measure (no fixed baseline source). A reduction theorem shows the
faithfulness constraints are free, so Hall's correlator-only threshold is
promoted to the faithful value; LP optima certified by integer arithmetic
give $F(n)=R/[2(R+1)]$ with $R=2^{\lfloor(n-1)/2\rfloor}$ (the {Mermin}
violation ratio) for $n=3,\dots,13$: $F(3)=1/3$, $F(5)=2/5$, $F(7)=4/9$,
$F(9)=8/17$, $F(11)=16/33$, $F(13)=32/65$ (each odd-$n$ value through $n=11$ repeated at
the following even size), with a universal ceiling $F\le 1/2$.

Our setup is different in two ways: the metric is per-block TV from a
fixed uniform source (not pairwise $L^{1}$), and the model class is plain
(no faithfulness constraints). The metrics are therefore
\emph{incommensurable} --- no ordering of the numbers follows. In
particular the numerical coincidence Alai $F(3)=1/3=$ our
$F^{*}(n=3;N=3)$ is a convention accident, not a theorem. What \emph{does}
align: both programs find that for GHZ-{Mermin} the minimum measurement
dependence \emph{does not vanish as the hidden source grows}, is determined
by the frustration structure of the inequality rather than the source
size, and admits a universal ceiling strictly below $1$ (theirs $\le 1/2$;
ours $\lim 3/8$ at $n=5$, oscillating cap at $n=3$). This paper
independently confirms, in a different metric, that ``exact minimum
measurement dependence for GHZ-{Mermin}'' is a well-posed problem with
clean structural answers --- and supplies the first TV-metric values for
it.

% =====================================================================
\section{Scope, caveats, and verification}
\label{sec:scope}

\subsection{Coverage limits, stated honestly}

\begin{itemize}
\item \textbf{{CGLMP} targets.} $T_{\mathrm{ME}}=(12+8\sqrt{3})/9$ is
  the maximally-entangled-state value of \cite{cglmp2002};
  $T_{\mathrm{Q}}=1+\sqrt{11/3}$ is the quantum maximum of \cite{adgl2002},
  matched to numerical precision by the NPA upper bound \cite{npa2008}; a
  closed-form optimality proof is not available, so if the true maximum
  $T'$ differed from $T_{\mathrm{Q}}$, then $F^{*}(T';N)=(T'-2)/(2\min(N,4))$ by the same
  linear-segment argument whenever $(T'-2)/8<\lfloor N/4\rfloor/N$ (for
  $N\ge 4$ this holds up to $T'<2+8\lfloor N/4\rfloor/N\ge 3$). The exact
  $\mathbb{Q}(\sqrt{3})$ certificate covers $T_{\mathrm{ME}}$; the
  $T_{\mathrm{Q}}$ values for $N=2,3$ rest on exhaustive floating-point
  enumeration. {CGLMP} $d\ge 4$ not treated.
\item \textbf{{Mermin} $n\ge 7$.} Pattern space $2^{n}$, rows $2^{n-1}$:
  full enumeration intractable. The structural claims ($m_{\min}$ grows
  with $n$; GHZ attains the algebraic maximum) were re-derived here only
  for $n=3,5$ (exact $\mathbb{Q}(i)$). Alai's $F(n)$ formula (Hall
  metric) covers $n\le 13$.
\item \textbf{Exact curves.} Certified at $N=2,\dots,6$ (all breakpoints;
  PM square to $N=10$). For $N\ge 7$: first segment and cap are proved
  for all $N$ (construction + global bound); the bend region is left to
  numerics (not computed here). The {Mermin} $n=5$ residues
  $r\in\{5,7,8,\dots,15\}\pmod{16}$ at large $m$ remain sandwiched between
  $\lceil 3N/8\rceil/N$ and $(6m+r)/(16m+r)$, not pinned.
\item \textbf{Model-class caveats.} Constructions use the uniform source;
  the lower-bound side does not depend on the source distribution (TV
  duality, same argument as \cite{mainpaper2026} --- stated, not
  re-certified per inequality). Stochastic responses reduce to the
  deterministic case by the vertex argument (same as
  \cite{mainpaper2026}) --- not separately certified here. The
  \emph{no-signalling sub-case} for these inequalities was \emph{not}
  studied (open; for {CHSH} it did not change $F^{*}$, so a similar
  result is plausible but unproved here).
\item \textbf{{I3322}} dropped (Sec.~\ref{sec:intro}); its local bound $0$
  was taken from the primary source's assertion, not re-proved by
  enumeration here.
\item \textbf{{KCBS} analog (open).} The PM square has quantum =
  algebraic maximum, so it tests only the cap. The {KCBS} inequality
  \cite{cabello2008} has a quantum value strictly between the NC bound
  and the algebraic maximum --- the true analog of $2\sqrt{2}$ vs $4$. The
  same two-lemma machinery applies (five contexts, binary outcomes,
  parity-type obstruction); the exact constants (NC bound of
  $\cos^{2}(\pi/10)$ type) must be taken from \cite{cabello2008} before
  use and were \emph{not} verified in this work.
\item \textbf{Novelty search.} {arXiv} API term combinations
  (``measurement dependence'' + ``contextuality''; ``fine-tuning'' +
  ``noncontextual'', date-sorted) found no prior work computing minimum
  measurement dependence / fine-tuning for these scenarios in the
  per-setting TV-from-a-fixed-baseline metric used here; the closest hit
  is \cite{alai2026b} (GHZ-{Mermin} Bell models, Hall's metric --- a
  different, incommensurable metric). Two term combinations only;
  paywalled/non-{arXiv} content not covered. This is a search result with
  stated coverage limits, not a proven novelty claim.
\end{itemize}

\subsection{Exact certificates (all re-runnable, all exit $0$)}

All results above are backed by re-runnable exact-arithmetic certificates
(exact $\mathbb{Q}$ or quadratic-field arithmetic; no floating point in
the proof parts; each script exits $0$ iff every assertion passes):

\footnotesize
\begin{equation*}
\begin{array}{l|l|l}
\text{Script} & \text{Log} & \text{Checks / runtime}\\
\hline
\texttt{ineq\_cglmp.py} & \texttt{ineq\_cglmp.out} & 30\ /\ \approx 25\ \mathrm{s}\\
\texttt{ineq\_mermin.py} & \texttt{ineq\_mermin.out} & 52\ /\ \approx 1\ \mathrm{s}\\
\texttt{ineq\_xcheck.py} & \texttt{ineq\_xcheck2.out} & 11\ /\ <1\ \mathrm{s}\\
\texttt{ineq\_mermin5\_n5.py} & \texttt{ineq\_mermin5\_n5.out} & 14\ /\ \approx 41\ \mathrm{s}\\
\texttt{fringe\_pm.py} & \texttt{fringe\_pm.out} & 88\ /\ \approx 1.6\ \mathrm{s}
\end{array}
\end{equation*}
\normalsize

\texttt{ineq\_cglmp.py}: difference-pattern reduction; full multiset
enumeration ($\mathbb{Q}$ + exact $\mathbb{Q}(\sqrt{3})$ tuples); per-
multiset $F^{*}$ analysis for $N=2,3$. \texttt{ineq\_mermin.py}:
$(P,r)$ pattern parameterization; full multiset enumeration; pins the
$n=5$ curves $N=2\dots 4$ and tightness of $\lceil 3N/8\rceil/N$;
enumerates $A_{\max B}$ exactly. \texttt{ineq\_xcheck.py}:
\emph{independent} cross-check by direct strategy brute force
($64/1024/81$ strategies, no pattern parameterization) and exact
$\mathbb{Q}(i)$ Kronecker computation of the GHZ correlators.
\texttt{ineq\_mermin5\_n5.py}: independent re-implementation of the
enumeration (cross-checked against the certified $N=4$ values before
use); full-curve pins $N=5,6$; audit of the remainder-witness
realizability. \texttt{fringe\_pm.py}: the PM square in full (exact
$\mathbb{Q}$ / $\mathbb{Q}(\sqrt{2})$; two soft LP cross-checks, never
failing the run). (The older log \texttt{ineq\_xcheck.out} is stale --- it
records a crash of an earlier buggy draft; the current output is
\texttt{ineq\_xcheck2.out}.) The two route families (pattern-
parameterized enumeration vs direct strategy brute force) agree on all
shared pattern-structure facts, and the GHZ quantum values were computed
without any shortcut formula. These certificates build on the main paper's
(the row lemma, the exact-curve breakpoint technique, the round-robin
witnesses), whose own certificates are listed in the appendices
of \cite{mainpaper2026}.

% =====================================================================
\section{Conclusion}
\label{sec:concl}

The minimum-fine-tuning program of \cite{mainpaper2026} generalizes
exactly where the quantum point lies on the first linear segment and the
obstruction has $m_{\min}=1$ balanced near-misses ({CGLMP} $d=3$), becomes
the cap $\lceil N/K\rceil/N$ where the quantum point \emph{is} the
algebraic maximum ({Mermin} $n=3$ with $K=4$; the Peres--Mermin square
with $K=6$, a first exact result for a contextuality scenario), and is
governed by the frustration index $m_{\min}/K$ where the obstruction is
frustrated ({Mermin} $n=5$, $\lim F^{*}=3/8$). The $K/m_{\min}$
principle (Conjecture~\ref{conj:kmm}) unifies the instances, with
{Mermin} $n=5$ marking where the finite-$N$ form can fail. Natural next
steps: the {KCBS} analog (the intermediate-quantum-value case for
contextuality), pinning the {Mermin} $n=5$ residues $r\in\{5,7,8,\dots,
15\}\pmod{16}$, and the no-signalling sub-case for these inequalities.

% =====================================================================
\section*{Note on provenance and corrections}

As explained in the corresponding note of the companion paper
\cite{mainpaper2026}, this paper is primarily the product of a locally
hosted open-weight language model (Qwen3.8-27B), which chose the problem,
did the mathematics and computations and wrote the text under the direction
of the human author, who supplied minimal scientific or technical input.
On 2026-09-03 it was critically reviewed by Claude Fable~5.1 (Anthropic),
which re-derived every result by independent exact computation (all
pattern censuses, the Mermin $n=5$ tables and $A_{\max B}$ values, the
Peres--Mermin curves for $N\le10$, the CGLMP closed forms) and found the
mathematics correct. The reviewer made the following corrections:
\begin{enumerate}
\item The CGLMP ``quantum value'' used as target,
  $T_{\mathrm{ME}}=(12+8\sqrt3)/9$, is the maximally-entangled-state value
  of Ref.~\onlinecite{cglmp2002}, not the quantum maximum: the larger value
  $T_{\mathrm{Q}}=1+\sqrt{11/3}$ of Ac\'in \emph{et al.}
  \cite{adgl2002} was missing (an earlier draft even recorded an
  unconfirmed guess ``$\approx3.34$''). Both targets are now treated; the
  closed form is unchanged, the $N\ge4$ headline number at $T_{\mathrm{Q}}$
  is $0.114357$.
\item Sec.~\ref{sec:setup} claimed that the grid values $V(k/N)$ always
  give the complete curve; this is false (the paper's own Mermin $n=5$,
  $N=2$ curve bends at $F=1/4$) and was replaced by the correct convexity
  statement.
\item Sec.~\ref{sec:cglmp}: the census omitted the six bound-attaining
  patterns of weight type $(+1,+1,0,0)$; now stated.
\item Theorem~\ref{thm:mermin5}: equality $F^{*}=\lceil 3N/8\rceil/N$ is
  provable for the residues $r\in\{0,1,2,3,4,6\}\pmod{16}$, not only
  $r\in\{0,1,2\}$; strengthened accordingly.
\item Sec.~\ref{sec:mermin3}: the combinatorial identity of the Mermin
  $n=3$ problem with CHSH is now stated.
\item References: the Peres--Mermin square is now credited to Peres (1990)
  and Mermin, PRL \textbf{65}, 3373, rather than to Mermin's GHZ paper; the
  Ac\'in \emph{et al.} reference was added; an over-generalisation of
  Alai's ``repeated at the following even size'' statement was restricted
  to $n\le11$ as in the source; two typographical errors
  ($\lfloor\cdot\rceil$, a missing $/N$) were fixed; the bibliography style
  was aligned with the companion paper.
\item The abstract was rewritten to state the contribution and its
  scope first, naming both {CGLMP} targets.
\end{enumerate}

\begin{acknowledgments}
This work received no institutional or financial support. As described in
the Note on provenance, the research and drafting were carried out by
language models under the author's direction; the author takes full
responsibility for the content.
\end{acknowledgments}

\bibliographystyle{apsrev4-2}
\bibliography{refs}

\end{document}
